
\subsubsection{6.1 Summary of Findings}

The experiments support the hypothesis that task-conditioned conversion can preserve adaptation capability at this scale:

\begin{enumerate}
\item Task structure is discoverable in pretrained hidden states (purity 0.74).
\item Contrastive learning produces embeddings with above-random discriminability (29.21\% vs 16.67\%).
\item Low-rank modulation achieves task-conditioned dynamics without overhead (0.85x).
\item TC-SSM achieves accuracy within 0.25\% of baseline and adapts in 14 steps.
\end{enumerate}

\subsubsection{6.2 Unexpected Findings}

The sub-unity overhead (0.85x) was not anticipated. Possible explanations include improved memory bandwidth utilization from batched task conditioning or measurement variance. Roofline analysis would be needed to confirm the mechanism.

\subsubsection{6.3 Limitations}

\textbf{Model scale.} All experiments use BERT-base (110M) to Mamba-130M conversion. Results may not generalize to larger models (1B+) where different phenomena may emerge.

\textbf{Task scope.} Evaluation covers SuperGLUE NLU classification only. Generation, reasoning, and multimodal tasks remain untested.

\textbf{Sample imbalance.} Small tasks (CB, COPA, WSC) achieve 0\% individual linear probe accuracy. Task embedding quality correlates with training data size. TC-SSM may underperform on rare task types.

\textbf{Baseline scope.} Comparison is against a transformer-style classifier baseline. Comparison against other conversion methods (e.g., LoRA post-conversion) was not conducted in the validation experiments, though the hypothesis synthesis document references Mamba + LoRA achieving 1.92\% gap.

\textbf{Architecture assumption.} Results are specific to Mamba-style SSM. Other sub-quadratic architectures (RWKV, linear attention) are not tested.

\subsubsection{6.4 Implications}

The results suggest that the efficiency-adaptation tradeoff in architecture conversion may not be fundamental, at least at this scale and for these task types. Task conditioning integrated during conversion preserves adaptation capability that post-hoc methods might struggle to recover.

